\documentclass[11pt]{article}
\usepackage{docstyle}
\newlength\figunit
\newcommand\ddxy{\dfrac{\mathrm{d}^{2}y}{\mathrm{d}x^{2}}}
\begin{document}
\sheethead{Lesson 2 \textperiodcentered\ Homework}{NCEA Level 2 \textperiodcentered\ AS 91262 \textperiodcentered\ Turning points and optimisation}{about 60 minutes}{}
Show all working. Justify every maximum or minimum. [A] Achieved, [M] Merit, [E] Excellence.

\Q Find the \tturn of $y=x^{2}+6x-7$ and state whether it is a maximum or a minimum.\mk{A}
\ALineT{$\dydx=2x+6=0$, so $x=-3$;\ \ $y=9-18-7=-16$}
\ALineT{$\ddxy=2>0$, so $(-3,\,-16)$ is a minimum}

\Q Find the \tturn of $y=3+10x-x^{2}$ and state its nature.\mk{A}
\ALineT{$\dydx=10-2x=0$, so $x=5$;\ \ $y=3+50-25=28$}
\ALineT{$\ddxy=-2<0$, so $(5,\,28)$ is a maximum}

\Q Find the \tturns of $f(x)=x^{3}-3x+2$ and justify their nature.\mk{A}
\ALine{$f'(x)=3x^{2}-3=0$, so $x=1$ or $x=-1$}
\ALine{$f(1)=1-3+2=0$;\quad $f(-1)=-1+3+2=4$}
\ALine{$f''(x)=6x$:\ $f''(1)=6>0$, so $(1,\,0)$ is a minimum}
\ALine{$f''(-1)=-6<0$, so $(-1,\,4)$ is a maximum}

\Q Find the \tturns of $f(x)=2x^{3}+3x^{2}-36x+1$ and justify their nature.\mk{A}
\ALine{$f'(x)=6x^{2}+6x-36=6(x+3)(x-2)=0$, so $x=-3$ or $x=2$}
\ALine{$f(-3)=-54+27+108+1=82$;\quad $f(2)=16+12-72+1=-43$}
\ALine{$f''(x)=12x+6$:\ $f''(-3)=-30<0$, so $(-3,\,82)$ is a maximum}
\ALine{$f''(2)=30>0$, so $(2,\,-43)$ is a minimum}

\Q Find the three \tturns of $y=x^{4}-8x^{2}$ and state their nature.\mk{A}
\ALineT{$\dydx=4x^{3}-16x=4x(x^{2}-4)=0$, so $x=0$, $x=2$ or $x=-2$}
\ALine{$y(0)=0$;\quad $y(\pm2)=16-32=-16$}
\ALineT{$\ddxy=12x^{2}-16$:\ at $x=0$ it is $-16<0$, so $(0,\,0)$ is a maximum}
\ALine{at $x=\pm2$ it is $32>0$, so $(2,\,-16)$ and $(-2,\,-16)$ are minimums}

\Q For which values of $x$ is $f(x)=x^{2}-4x+9$ decreasing?\mk{A}
\ALine{$f'(x)=2x-4<0$, so $x<2$}
\ALine{}

\Q For which values of $x$ is $f(x)=x^{3}-3x^{2}+5$ increasing?\mk{A}
\ALine{$f'(x)=3x^{2}-6x=3x(x-2)$, zero at $x=0$ and $x=2$}
\ALine{Up-parabola, positive outside its roots:\ $x<0$ or $x>2$}

\Q For which values of $x$ is $f(x)=\dfrac{1}{3}x^{3}-x^{2}-8x$ decreasing?\mk{A}
\ALine{$f'(x)=x^{2}-2x-8=(x-4)(x+2)$, zero at $x=4$ and $x=-2$}
\ALine{Negative between the roots:\ $-2<x<4$}

\Q The curve $y=x^{2}+bx+5$ has a \tturn at $x=3$. Find $b$ and the $y$-coordinate of the \tturn.\mk{M}
\ALineT{$\dydx=2x+b=0$ at $x=3$:\ $6+b=0$, so $b=-6$}
\ALine{$y=9-18+5=-4$;\ the \tturn is $(3,\,-4)$}

\Q A model rocket's height is $h=40t-5t^{2}$ metres after $t$ seconds. Find its maximum height.\mk{A}
\ALineT{$\dd{h}{t}=40-10t=0$, so $t=4$ s}
\ALine{$h=160-80=80$;\ $h''=-10<0$, so maximum height $80$ m}

\Q The daily cost of making $x$ units is $C=2x^{2}-80x+1000$ dollars. Find the number of units that gives the minimum cost, and the minimum cost.\mk{A}
\ALineT{$\dd{C}{x}=4x-80=0$, so $x=20$ units}
\ALine{$C=800-1600+1000=200$;\ $C''=4>0$, so the minimum cost is \$200}

\Q Two positive numbers add to $20$. Use calculus to find the largest possible value of their product.\mk{M}
\ALine{Numbers $x$ and $20-x$:\ $P=x(20-x)=20x-x^{2}$}
\ALine{$P'=20-2x=0$, so $x=10$;\ $P''=-2<0$, so maximum}
\ALine{Largest product $=10\times10=100$}

\Q A rectangular garden against a house wall is fenced on the other three sides with $40$ m of fencing. Find the dimensions that give the largest area, and that area.\mk{M}
\ALine{Sides $x$, $x$ and $40-2x$:\ $A=x(40-2x)=40x-2x^{2}$}
\ALine{$A'=40-4x=0$, so $x=10$;\ $A''=-4<0$, so maximum}
\ALine{Garden $10$ m by $20$ m; area $200\ \mathrm{m}^{2}$}

\QX A rectangle is drawn under the curve $y=12-x^{2}$ with two corners on the $x$-axis and two corners on the curve, symmetrical about the $y$-axis. Find the largest possible area of the rectangle, and justify that it is a maximum.\mk{E}

\begin{minipage}[t]{0.54\linewidth}
\ALine{Corner $(x,\ 12-x^{2})$:\ width $2x$, height $12-x^{2}$}
\ALine{$A=2x(12-x^{2})=24x-2x^{3}$}
\ALine{$A'=24-6x^{2}=0$, so $x^{2}=4$, $x=2$ ($x>0$)}
\ALine{$A''=-12x=-24<0$, so maximum}
\ALine{Largest area $=2(2)(12-4)=32$ square units}
\end{minipage}\hfill
\begin{minipage}[t]{0.44\linewidth}\centering\vspace{0pt}
\setlength\figunit{7mm}
% Rectangle under y = 12 - x^2, vertices (+-a, 0) and (+-a, 12 - a^2), drawn with a = 1.6.
\begin{tikzpicture}[x=\figunit, y=0.25\figunit, line cap=round]
  \draw[-{Stealth[length=2mm]}, line width=0.6pt] (-4.2,0) -- (4.4,0) node[right, inner sep=1pt] {$x$};
  \draw[-{Stealth[length=2mm]}, line width=0.6pt] (0,-1) -- (0,14) node[above, inner sep=1pt] {$y$};
  \begin{scope}\clip (-4.2,-1) rectangle (4.4,13.5);
    \curve{-3.8}{3.8}{12-\x*\x}\end{scope}
  \def\a{1.6}
  \draw[line width=0.9pt, fill=boxfill] (-\a,0) rectangle (\a,{12-\a*\a});
  \fill (\a,{12-\a*\a}) circle (2pt);
  \node[right, inner sep=3pt] at (\a,{12-\a*\a}) {$(x,\ 12-x^{2})$};
  \node[anchor=south west] at (2.9,{12-2.9*2.9}) {$y=12-x^{2}$};
\end{tikzpicture}

\end{minipage}

\vfill
{\small Optional extra: StudyTime Level 2 Calculus Guide, ``Calculating Turning Points'' from page 16.\par}
\end{document}
